\section{Containerization}

\subsection{Building the Backend Container}

The backend needed a multi-stage Containerfile. The reason for two stages is practical:
the first stage uses a large Go toolset image to compile the code, and the second stage
takes only the resulting binary and puts it into a minimal Red Hat UBI image. The final
image ends up around 94 MB instead of over a gigabyte, which matters for pull times and
attack surface.

In the build stage I used the Red Hat Golang toolset image and followed the pattern of
copying \texttt{go.mod} and \texttt{go.sum} first, running \texttt{go mod download}, and
only then copying the source. This way Docker can cache the dependency layer and skips
the download step on rebuilds unless the module files change. The binary was built with
\texttt{CGO\_ENABLED=0} and \texttt{GOOS=linux} so it comes out as a fully static
binary with no runtime dependencies. The \texttt{-ldflags} and \texttt{-trimpath} flags
strip debug paths from the binary, which is cleaner for production.

\subsection{Architecture and Permission Challenges}

Several problems came up during the build that were not obvious from the instructions.

The first was a missing lockfile. When I ran the build for the first time it failed with
\texttt{copier: stat: "/go.sum": no such file or directory}. The repository was new and
\texttt{go.sum} had never been generated. Running \texttt{go mod tidy} in the
\texttt{backend/} directory fixed that. It is a small thing but easy to miss if you are
not used to Go projects.

The second problem was a permission error during the build phase. The error was
\texttt{open backend: permission denied}. The Red Hat UBI images run as a non-root user
by default, UID 1001, and that user cannot write to \texttt{/app} which is owned by root.
The fix was to change \texttt{WORKDIR} to \texttt{/opt/app-root/src}, which is the
directory that the UBI Go toolset image already prepares for exactly this purpose. I had
not known that convention before, but it makes sense once you understand that Red Hat
designs its images to be secure out of the box rather than trusting the developer to add
security later.

The third issue was an architecture mismatch. My machine is an ARM64 Mac, but the
original Containerfile had \texttt{GOARCH=amd64} hardcoded. When I ran the resulting
container I got an \texttt{Illegal Instruction} crash immediately. The binary had been
compiled for a CPU that the machine did not have. Removing the hardcoded \texttt{GOARCH}
lets Go detect the host architecture automatically and build a native ARM64 binary, which
runs without emulation and is noticeably faster.

\subsection{Frontend Container}

The frontend container uses an Nginx base image from Red Hat. One thing I ran into was
that the container started and then immediately exited with status 0. The reason is that
the Red Hat Nginx image is designed as a Source-to-Image builder, so without an explicit
\texttt{CMD} it does not keep the daemon running in the foreground. Adding the following
line to the Containerfile solved it:

\begin{lstlisting}[language=bash]
CMD ["nginx", "-g", "daemon off;"]
\end{lstlisting}

I also needed to make sure the \texttt{nginx.conf} was copied to
\texttt{/etc/nginx/nginx.conf} so Nginx would pick it up.

\subsection{Nginx as a Reverse Proxy}

Nginx does more than serve the static frontend files. It also proxies API calls from the
browser to the Go backend. Without this, the browser would try to call
\texttt{/api/entries} directly on the same host as the frontend, and there would be
nothing there to answer. The relevant part of \texttt{nginx.conf} looks like this:

\begin{lstlisting}
location /api/ {
    proxy_pass http://backend;
    proxy_set_header Host $host;
}

location /health {
    proxy_pass http://backend;
}
\end{lstlisting}

The hostname \texttt{backend} works because all containers share the same network.
Locally that network is created with \texttt{podman network create}, and on OpenShift it
is handled automatically by the cluster's internal DNS.

\subsection{Local Testing with Podman}

Before moving to OpenShift I tested the full stack locally using Podman. The following
commands set everything up from scratch:

\begin{lstlisting}[language=bash]
# Create shared network and data directories
podman network create guestbook-net || true
mkdir -p ~/guestbook-data/pgdata ~/guestbook-data/redisdata

# Build both application images
podman build -t guestbook-backend ./backend
podman build -t guestbook-frontend ./frontend

# PostgreSQL
podman run -d --name postgres \
  --network guestbook-net \
  -e POSTGRESQL_USER=guestbook \
  -e POSTGRESQL_PASSWORD=password \
  -e POSTGRESQL_DATABASE=guestbook \
  -v ~/guestbook-data/pgdata:/var/lib/pgsql/data:Z \
  quay.io/fedora/postgresql-16:latest

# Redis
podman run -d --name redis \
  --network guestbook-net \
  -e REDIS_PASSWORD=password \
  quay.io/kurs/redis:latest

# Backend
podman run -d --name backend \
  --network guestbook-net \
  -e DB_HOST=postgres \
  -e DB_USER=guestbook \
  -e DB_PASSWORD=password \
  -e REDIS_HOST=redis \
  -e REDIS_PASSWORD=password \
  localhost/guestbook-backend

# Frontend
podman run -d --name frontend \
  --network guestbook-net \
  -p 8080:8080 \
  localhost/guestbook-frontend
\end{lstlisting}

One detail worth noting is the \texttt{:Z} flag on the PostgreSQL volume. On systems
running SELinux, which includes Fedora and RHEL and therefore the nodes that OpenShift
runs on, the kernel prevents containers from writing to host directories by default.
The \texttt{:Z} flag tells Podman to relabel the directory with the correct SELinux
context. Without it, PostgreSQL crashes with a permission error when it tries to
initialize its data directory, regardless of which user you are running as.

Another early mistake was forgetting to add \texttt{-p 8080:8080} to the frontend run
command. The container was running fine internally but nothing was reachable from the
browser. The fix was straightforward once I understood that port mapping is not automatic.

\subsection{Cache Invalidation in the Backend}

The backend implements a read-through cache. When the frontend requests the list of
entries, the backend checks Redis first. If the data is there it returns it immediately
without touching PostgreSQL. If not, it queries PostgreSQL and then stores the result
in Redis so the next request is fast. The important counterpart to this is what happens
when someone adds a new entry. If the cached list is not deleted at that point, users
will see stale data until the cache expires naturally. The relevant code in
\texttt{main.go} handles this:

\begin{lstlisting}[language=Go]
// Invalidate cache on write
if app.Redis != nil {
    app.Redis.Del(app.Ctx, "entries:all")
}
\end{lstlisting}

This is a classic cache invalidation pattern and something I had read about but not
implemented before. Seeing it break in practice when I accidentally left it out made the
concept stick much more clearly.
